\section{Transfer of the first unsolved word coefficient}
\label{sec:word-transfer}

Our next objective is to control equations in free letters.  The first and
paired responses describe matrix perturbations with one loop or two opposite
arrows.  We show that they also contain the first unsolved coefficient of an
arbitrary free-series perturbation.  The argument unfolds a word along a
triangular chain and then replaces its internal states by a single auxiliary
partner.  The only difficulty is that this partner initially has the target
eigenvalue.  The following linear-algebra observation explains how to remove
that singularity without changing the coefficient being computed.

\begin{lemma}[An inverse between two vanishing couplings]
\label{lem:killed-pole}
Let $E$ be a field and let $M(\xi)\in\operatorname{End}_E(V)[[\xi]]$, where
$V$ is finite dimensional.  Suppose that zero is a semisimple eigenvalue of
$M_0=M(0)$, and let $\Pi$ be its spectral projection.  Suppose that
\[
 G=\Pi M'(0)\Pi\big|_{\ker M_0}
\]
is invertible.  Then $M(\xi)^{-1}$ exists over $E((\xi))$ and has a pole of
order at most one.  Let $J$ be zero on $\ker M_0$ and equal to $M_0^{-1}$ on
$\operatorname{im}M_0$.  If $L(\xi)$ and $B(\xi)$ are compatible matrix
series satisfying
\[
 L(0)\Pi=0,\qquad \Pi B(0)=0,
\]
then
\begin{equation}
\label{eq:killed-pole-limit}
 \lim_{\xi\to0}L(\xi)M(\xi)^{-1}B(\xi)=L(0)JB(0).
\end{equation}

Suppose additionally that $e$ is an idempotent, $f=1-e$,
$fM_0f$ is invertible on $fV$, and $e:\ker M_0\to eV$ is an isomorphism.
Under the same two coupling conditions,
\begin{equation}
\label{eq:fixed-complement-inverse}
 L(0)JB(0)=L(0)f(fM_0f)^{-1}fB(0).
\end{equation}
\end{lemma}

\begin{proof}
Use $V=\ker M_0\oplus\operatorname{im}M_0$.  In this decomposition
\[
 M(\xi)=
 \begin{pmatrix}
  \xi G+O(\xi^2)&O(\xi)\\
  O(\xi)&M_0|_{\operatorname{im}M_0}+O(\xi)
 \end{pmatrix}.
\]
The lower-right block is invertible and the Schur complement of that block
is $\xi G+O(\xi^2)$.  The upper-left block of $M(\xi)^{-1}$ is consequently
$\xi^{-1}G^{-1}+O(1)$; its off-diagonal blocks are bounded; and its
lower-right block tends to $M_0^{-1}|_{\operatorname{im}M_0}$.
The kernel components of both couplings are $O(\xi)$.  Multiplying these
four blocks proves \eqref{eq:killed-pole-limit}.

For the second assertion, write $b=B(0)u$.  The condition $\Pi b=0$ places
$b$ in $\operatorname{im}M_0$.  Solutions of $M_0x=b$ differ by
$\ker M_0$.  Since projection to $eV$ is an isomorphism on that kernel,
there is a unique solution with $ex=0$.  Its $f$ coordinate is
$(fM_0f)^{-1}fb$.  All solutions have the same image under $L(0)$ because
$L(0)$ kills $\ker M_0$.  Comparing this solution with $Jb$ proves
\eqref{eq:fixed-complement-inverse}.
\end{proof}

We now recall exactly which response identities enter the transfer
argument.  At a small generic point $P$ of size $n\in\mathcal S$, the
eigenvalues $\lambda_h$ have projections $\pi_h$ of rank $t$.  The first
response is $D_h(P;Y)\pi_h$.  A pair of arrows through an independent
partner $Q$ gives $m_h(P,Q;U,V)\pi_h$.  The columns of the first responses
and the coefficient vectors of all these paired responses span
$\mathcal T_P$.  By Lemma~\ref{lem:active-specialization}, their scalarity
and all linear containment identities persist for a block-triangular
partner after its finitely many resolvent determinants are cleared.  For a
set of target labels $I$, only the determinants belonging to labels in $I$
are needed.  We use the notation
\[
 \mathcal T_P|_I=\{(c_h)_{h\in I}:(c_h)_h\in\mathcal T_P\}.
\]
These spaces and identities are extended to the coefficient field in use.
We also fix a direction $v$ with $D_h(P;v)\ne0$ for every target label.

\begin{proposition}[Transfer of a leading coefficient]
\label{prop:word-transfer}
Let $R=\operatorname{diag}(P,C_0)$ be a small matrix input of total size $N\in
\mathcal S$;
the block $C_0$ may be absent.  Let $I$ be a set of labels at
$P$ such that $W(C_0)-\lambda_h$ is invertible for every $h\in I$.
Work over a field $E$ containing the labeled spectral data at $P$ and the
entries of $C_0$, with the formal evaluations allowed by
Lemma~\ref{lem:admissible-evaluation}.  Let
\[
 Z(s_1,s_2)\in M_N(E)\langle\!\langle s_1,s_2\rangle\!\rangle^2
\]
have constant term $R$.  These are input matrices: their scalar entries
act on the input factor and commute with the old coefficient algebra.

For $h\in I$, embed $\pi_h$ at the $P$ block, put $\bar\pi_h=1-\pi_h$,
and form
\[
 S_h(Z)=\pi_h(W(Z)-\lambda_h)\pi_h
 -\pi_hW(Z)\bar\pi_h
 \bigl(\bar\pi_h(W(Z)-\lambda_h)\bar\pi_h\bigr)^{-1}
 \bar\pi_hW(Z)\pi_h.
\]
Suppose that every $S_h(Z)$ vanishes in degrees less than $\ell$, where
$\ell\ge1$.  For each word $\omega$ of length $\ell$, there are scalars
$c_{h,\omega}\in E$ such that
\[
 [S_h(Z)]_\omega=c_{h,\omega}\pi_h,
 \qquad (c_{h,\omega})_{h\in I}\in
        (\mathcal T_P|_I)\otimes E.
\]
The same conclusion holds for homogeneous degree-$\ell$ coefficients of
perturbations in finitely many ordinary commuting parameters.
\end{proposition}

\begin{proof}
For $\ell=1$, only the $P$-to-$P$ input coefficient contributes to the
compressed derivative, so the assertion is the first-response identity.
Assume $\ell\ge2$ and fix a word $\omega$ of length $\ell$.

\paragraph{Unfolding the word.}
Use $\ell+1$ ordered states.  On consecutive steps put the letters of
$\omega$, and evaluate a word by following that path.  Thus a coefficient
of $Z$ for a substring becomes the corresponding strictly upper block of
the input, while every diagonal block is $R$.  The $(1,\ell+1)$ block
extracts precisely the coefficient of $\omega$.  This substitution is a
homomorphism on the free-series ring modulo degree $\ell+1$, and it
therefore respects every inverse with invertible constant term.

Keep only $P$ at the two endpoint states, leaving all $R$ blocks at the
$\ell-1$ internal states.  This operation does not change the required
compressed endpoint coefficient.  At an endpoint, constant input steps
preserve the blocks $P,C_0$, and the constant output commutes with
$\pi_h$.  A contributing Schur path can therefore neither enter the
first endpoint complement from its target projection nor leave the last
endpoint complement for its target projection.  All complement
eliminations relevant to the coefficient occur at internal states.

Denote the internal input by $Q_0$.  Its size is $(\ell-1)N$ and hence
belongs to $\mathcal S$.  Grouping the internal
states, write the cropped input as
\[
 T_0=\begin{pmatrix}
 P&G_1&Y\\0&Q_0&G_2\\0&0&P
 \end{pmatrix}.
\]
Here each entry denotes a pair of input matrices.  In the output write
$E_1=W(T_0)_{12}$, $E_2=W(T_0)_{23}$, and $F=W(T_0)_{13}$.

\paragraph{The internal target eigenspace.}
Fix $h\in I$, set $M_0=W(Q_0)-\lambda_h$, and let $e$ be the sum of the
fixed copies of $\pi_h$ on its internal states.  Put $f=1-e$.
The matrix $fM_0f$ is upper triangular with invertible diagonal blocks.
Moreover its Schur complement on $eV$ is zero: every strictly upper
coefficient is a coefficient for a proper substring of $\omega$, and
these vanish by hypothesis.  Its kernel is consequently the graph
\[
 u\longmapsto u-f(fM_0f)^{-1}fM_0e u,
 \qquad u\in eV.
\]
Projection onto $eV$ is an isomorphism on this kernel, whose dimension is
$(\ell-1)t$.  Triangularity shows that this is also the full algebraic
multiplicity of zero.  Thus zero is semisimple.  Let $\Pi$ be its
spectral projection.

The same argument applies to the chain consisting of the first endpoint
and the internal states, and to the chain consisting of the internal
states and the last endpoint.  They involve only proper substrings.
Solving their triangular kernel equations shows that the couplings
between their zero eigenspaces vanish:
\begin{equation}
\label{eq:transfer-coupling-annihilation}
 \pi_hE_1\Pi=0,\qquad \Pi E_2\pi_h=0.
\end{equation}
Indeed any nonzero compressed coupling would reduce the kernel dimension
below the sum of the diagonal zero multiplicities.  This argument uses
semisimplicity only at the target eigenvalue.

\paragraph{A regularized auxiliary partner.}
In every internal copy of $R$, replace its $P$ block by $P+\xi v$, keeping
all arrows fixed; call the resulting input $Q(\xi)$.  Put
$M(\xi)=W(Q(\xi))-\lambda_h$.
The spectral projection $\Pi$ is upper triangular in the state order,
with diagonal equal to the fixed target projections.  The graph
identification of its range with $eV$ is upper triangular with identity
diagonal.  It follows that
\[
 \Pi M'(0)\Pi\big|_{\ker M_0}
\]
is upper triangular in this identification, with every diagonal block
equal to $D_h(P;v)I_t$.  It is therefore invertible.  This verifies the
simple-pole hypothesis of Lemma~\ref{lem:killed-pole}, including the
effect of all internal triangular arrows.

Let $T(\xi)$ be the corresponding full cropped input and use $E_1(\xi)$,
$E_2(\xi)$, $F(\xi)$ for its output blocks.  For nonzero formal $\xi$, the
compressed endpoint expression is
\begin{align}
 \mathcal E_h(\xi)
 &=\pi_h\bigl(F(\xi)-E_1(\xi)M(\xi)^{-1}E_2(\xi)\bigr)\pi_h
 \label{eq:transfer-regularized-expression}\\
 &=D_h(P;Y)\pi_h+M_h(P,Q(\xi);G_1,G_2).
 \nonumber
\end{align}
Its tuple is scalar and belongs to the projected active space over
$E((\xi))$.  This is exactly the specialized paired-response identity:
the auxiliary block is upper triangular, its diagonal entries are small
blocks with a Taylor shift, and its strictly upper entries are fixed
finite matrices.  Clear the target resolvent determinants first and
apply Lemma~\ref{lem:active-specialization}; restore the inverses over
$E((\xi))$, where the preceding argument proves they exist.  No
specialization of an undefined inverse is being used.

\paragraph{Identifying the limit.}
By \eqref{eq:transfer-coupling-annihilation} and
Lemma~\ref{lem:killed-pole}, expression
\eqref{eq:transfer-regularized-expression} has the finite limit
\[
 \pi_h\bigl(F-E_1JE_2\bigr)\pi_h
 =\pi_h\bigl(F-E_1f(fM_0f)^{-1}fE_2\bigr)\pi_h.
\]
The last expression eliminates precisely the fixed complementary
spaces in the internal states.  By the unfolding and endpoint-cropping
argument it is $[S_h(Z)]_\omega$.  In particular the limiting inverse has
not changed the coefficient to one computed with a different choice of
complement.  Taking the limit in the scalar-response identities and in
the finitely many linear equations defining $\mathcal T_P|_I$ proves
the assertion simultaneously for all target labels.

Finally consider a perturbation in commuting variables $h_1,\ldots,h_s$.
Substitute $h_j=a_jz$ for arbitrary scalar directions $a_j$ and one free
letter $z$.  The lower homogeneous terms still vanish, so the result
just proved applies to the degree-$\ell$ coefficient on every such
line.  Scalarity and containment in a fixed finite-dimensional linear
space are linear conditions.  Each of their annihilators therefore
gives a homogeneous polynomial in the $a_j$ that vanishes identically.
Over the infinite characteristic-zero field all its coefficients vanish.
This proves the commuting-parameter assertion.
\end{proof}

Proposition~\ref{prop:word-transfer} is the bridge from responses to
free-series equations.  It controls a first unsolved coefficient without
assuming that the perturbation is commutative or that the auxiliary
matrix has a separated spectrum before regularization.  We next use its
commuting-parameter consequence to show that linked responses arise from
actual formal relations between eigenvalues.
